\documentclass[1pt]{article}
\usepackage[a4paper,margin=0.2in]{geometry} % Small margins
\usepackage{multicol} % For multiple columns
\usepackage{amsmath, amssymb} % For math symbols
\usepackage{bm} % Bold math
\usepackage{graphicx} % To insert small images
\usepackage{array}
\usepackage{booktabs}
\usepackage{placeins}
\usepackage{enumitem} % To make lists compact
\setlength{\columnsep}{1pt} % Reduce space between columns
\raggedcolumns{}
\begin{document}
\begin{multicols}{3} % 3 columns
\raggedcolumns
\setlength{\columnseprule}{0.1pt}
\scriptsize % Use a smaller font
\noindent % Avoid extra space at beginning
\abovedisplayskip=2pt
\belowdisplayskip=2pt

%====================================================================
% I. Mechanics Fundamentals
%====================================================================

\textbf{Order of Operations}
\begin{enumerate}[noitemsep, topsep=0pt]
    \item \textbf{Counting Degrees of Freedom:}
    holonomic constraint:
    \begin{equation*}
        F\left(\vec{r}_1, \vec{r}_2, \ldots, \vec{r}_N, t\right)=0
    \end{equation*}
    number of free DOF \( n=3N-k \) where \( k \) is the number of constraints.
    \item Find \( n \) generalized coordinates.  
    \item Express in terms of generalized coordinates.  
    \item Calculating the Kinetic and Potential Energy.  
    \item Writing the Equations of Motion.
\end{enumerate}

\textbf{Generalized Kinetic Energy}
\begin{align*}
     T = \frac{1}{2} m (\dot{x}^2 + \dot{y}^2 + \dot{z}^2) \quad \text{(Cartesian)} 
\end{align*}
\begin{align*}
    T = \frac{1}{2} m (\dot{r}^2 + r^2 \dot{\theta}^2 + r^2 \sin^2 \theta \dot{\phi}^2) \quad \text{(Spherical)}
\end{align*}
\begin{align*}
    T = \frac{1}{2} m (\dot{r}^2 + r^2 \dot{\theta}^2 + \dot{z}^2) \quad \text{(Cylindrical)}
\end{align*}


\textbf{(M and K Matrices)}
\begin{equation*}
       T = \frac{1}{2} M_{ij} \dot{q}_i \dot{q}_j = \frac{1}{2} \mathbf{\dot{q}}^T \mathbf{M} \mathbf{\dot{q}}
\end{equation*}
\begin{equation*}
    V = \frac{1}{2} K_{ij} q_i q_j = \frac{1}{2} \mathbf{q}^T \mathbf{K} \mathbf{q}
\end{equation*}
\begin{equation*}
M_{ij} = \frac{\partial^2 \mathcal{L}}{\partial \dot{q}_i \partial \dot{q}_j} \quad
K_{ij} = -\frac{\partial^2 \mathcal{L}}{\partial q_i \partial q_j} 
\end{equation*}
\[
M_{ij} =
\begin{bmatrix}\tiny{}

\frac{\partial^2 \mathcal{L}}{\partial \dot{q}_1 \partial \dot{q}_1} & \frac{\partial^2 \mathcal{L}}{\partial \dot{q}_1 \partial \dot{q}_2} & \frac{\partial^2 \mathcal{L}}{\partial \dot{q}_1 \partial \dot{q}_3} \\
\frac{\partial^2 \mathcal{L}}{\partial \dot{q}_2 \partial \dot{q}_1} & \frac{\partial^2 \mathcal{L}}{\partial \dot{q}_2 \partial \dot{q}_2} & \frac{\partial^2 \mathcal{L}}{\partial \dot{q}_2 \partial \dot{q}_3} \\
\frac{\partial^2 \mathcal{L}}{\partial \dot{q}_3 \partial \dot{q}_1} & \frac{\partial^2 \mathcal{L}}{\partial \dot{q}_3 \partial \dot{q}_2} & \frac{\partial^2 \mathcal{L}}{\partial \dot{q}_3 \partial \dot{q}_3}
\end{bmatrix}
\]
\[
K_{ij} =
-\begin{bmatrix}
\frac{\partial^2 \mathcal{L}}{\partial q_1 \partial q_1} & \frac{\partial^2 \mathcal{L}}{\partial q_1 \partial q_2} & \frac{\partial^2 \mathcal{L}}{\partial q_1 \partial q_3} \\
\frac{\partial^2 \mathcal{L}}{\partial q_2 \partial q_1} & \frac{\partial^2 \mathcal{L}}{\partial q_2 \partial q_2} & \frac{\partial^2 \mathcal{L}}{\partial q_2 \partial q_3} \\
\frac{\partial^2 \mathcal{L}}{\partial q_3 \partial q_1} & \frac{\partial^2 \mathcal{L}}{\partial q_3 \partial q_2} & \frac{\partial^2 \mathcal{L}}{\partial q_3 \partial q_3}
\end{bmatrix}
\]
\begin{align*}
    M_{ij} \ddot{q}_j + K_{ij} q_j = 0
\end{align*}
\textbf{Equilibrium and Potential Approximation}
\begin{enumerate}
    \item \textbf{Find equilibrium:} Solve $\nabla V(\mathbf{r}) = 0$.
    \item \textbf{Classify equilibrium:} Compute Hessian $K = \nabla^2 V(\mathbf{r})$.
    \begin{itemize}
        \item Positive eigenvalues $\Rightarrow$ Stable (minimum).
        \item Negative eigenvalues $\Rightarrow$ Unstable (maximum).
        \item Mixed eigenvalues $\Rightarrow$ Saddle (unstable).
    \end{itemize}
    \item \textbf{Approximate potential:} Use Taylor expansion near $\mathbf{r}_0$:
    \[
    V(\mathbf{r}) \approx V(\mathbf{r}_0) + \frac{1}{2} (\mathbf{r} - \mathbf{r}_0)^T K (\mathbf{r} - \mathbf{r}_0).
    \]
    \item \textbf{Small oscillations:} Solve $M \ddot{\mathbf{r}} = -K(\mathbf{r} - \mathbf{r}_0)$ to find normal modes.
\end{enumerate}

\textbf{Generalized Velocity and Acceleration}
\[\dot{\mathbf{r}}
\;=\;
\frac{\partial \mathbf{r}}{\partial \theta}\,\dot{\theta}
\;+\;
\frac{\partial \mathbf{r}}{\partial \phi}\,\dot{\phi},\quad     \dot{r}_a = \sum_i \frac{\partial r_a}{\partial q_i} \dot{q}_i
\]
\begin{equation*}
\end{equation*}
\begin{equation*}
    \ddot{r}_a = \sum_i \frac{\partial r_a}{\partial q_i} \ddot{q}_i + \sum_{i,j} \frac{\partial^2 r_a}{\partial q_i \partial q_j} \dot{q}_i \dot{q}_j
\end{equation*}

\textbf{Generalized Force}
\begin{equation*}
    Q_{j}=\sum_{i}\vec{F}_{i}\frac{\partial r_{i}}{\partial q_{j}}=-\frac{\partial V}{\partial q_j}
\end{equation*}
\begin{equation*}
    Q_j^{(nc)}=\sum_i \vec{F}_i^{(nc)} \cdot \frac{\partial \vec{v}_i}{\partial \dot{q}_j}
\end{equation*}

\textbf{Conservative and Central Forces} 
\(  \vec{J}=\intop_{t_{1}}^{t_{2}}F\,dt\)  \quad

\( W = \int \mathbf{F} \cdot d\mathbf{r} \) \quad  \quad
\[
\mathbf{F} = -\nabla V
\]
\[\delta V=\int\delta Fdr,\quad V=\int dV=\iint\delta Fdr\]

\textbf{Angular Momentum and Torque}
\begin{equation*}
    \mathbf{L} = \mathbf{r} \times \mathbf{p} = mr^2 = I\omega
\end{equation*}
\begin{equation*}
    \tau= \mathbf{r} \times \mathbf{F} = \frac{d\mathbf{L}}{dt} = mr^{2}\dot{\omega} = \mathbf{I}\dot{\omega}
\end{equation*}
\[L_{rot}=I\vec{\omega}=\begin{pmatrix}I_{x}\omega_{x}\\
I_{y}\omega_{y}\\
I_{z}\omega_{z}
\end{pmatrix}\]

\textbf{Non-Inertial Systems}

\[
\ddot{\mathbf{r}} = 
\underbrace{m\ddot{\mathbf{r}}'}_{\mathbf{F}_{\text{ext}}} 
 \underbrace{ -2m\boldsymbol{\omega} \times \dot{\mathbf{r}}'}_{\mathbf{F}_{\text{cor}}} 
 \underbrace{ -m\dot{\boldsymbol{\omega}} \times \mathbf{r}'}_{\mathbf{F}_{\text{eul}}} 
 \underbrace{ -m\boldsymbol{\omega} \times (\boldsymbol{\omega} \times \mathbf{r'})}_{\mathbf{F}_{\text{cent}}}
\]
\[\mathcal{L}_{rot}=\frac{m}{2}\left(\vec{v}+\vec{\Omega}\times\vec{r}\right)^{2}-V\left(\vec{r}\right)\]
\[\mathcal{H}_{rot}=\frac{p^{2}}{2m}-\vec{p}\cdot\left(\vec{\Omega}\times\vec{r}\right)+V\left(\vec{r}\right)\]

%====================================================================
% II. Lagrangian and Hamiltonian Mechanics
%====================================================================

\textbf{Lagrangian Mechanics}
\begin{align*}
    \frac{\partial \mathcal{L}}{\partial q} - \frac{d}{dt}\left(\frac{\partial \mathcal{L}}{\partial \dot{q}}\right) + \frac{d^{2}}{dt^{2}}\left(\frac{\partial \mathcal{L}}{\partial \ddot{q}}\right)=0
\end{align*}

\textbf{With Force}
\begin{equation*}
    \frac{d}{dt}\left[\frac{\partial T}{\partial\dot{q}_{j}}\right]-\frac{\partial T}{\partial q_{j}}=-\frac{\partial V}{\partial q_j}=Q_{j}
\end{equation*}
\begin{equation*}
     \frac{d}{dt}\left[\frac{\partial \mathcal{L}}{\partial \dot{q}_j}\right]-\frac{\partial \mathcal{L}}{\partial q_j}=Q_j^{(nc)}=\sum_{i}\vec{F}_{i}\frac{\partial r_{i}}{\partial q_{j}}
\end{equation*}

\textbf{Rayleigh Function and Dissipation}
\begin{equation*}
    F = \frac{1}{2} \sum_{i,j} C_{ij} \dot{q}_i \dot{q}_j \quad \text{or} \quad F = \frac{1}{2} \sum_{i} C_i \dot{q}_i^2
\end{equation*}
\begin{equation*}
    Q_i^{\text{(diss)}} = - \frac{\partial F}{\partial \dot{q}_i} = - c \dot{x}
\end{equation*}
\begin{equation*}
    \frac{d}{dt} \left( \frac{\partial \mathcal{L}}{\partial \dot{q}_i} \right) - \frac{\partial \mathcal{L}}{\partial q_i} = -\frac{\partial F}{\partial \dot{q}_i}
\end{equation*}
\[\frac{d}{dt}\left(\frac{\partial\mathcal{L}}{\partial\dot{q}_{i}}\right)-\frac{\partial\mathcal{L}}{\partial q_{i}}+\frac{\partial\mathcal{R}}{\partial\dot{q_{i}}}=F\]
\begin{equation*}
    P_{\text{diss}} = \sum_{i} Q_i^{\text{(diss)}} \dot{q}_i = 2F
\end{equation*}

\textbf{Cyclic Coordinates, Canonical Momentum}
\begin{equation*}
    \frac{\partial \mathcal{L}}{\partial q_i} = 0 \implies \frac{d}{dt}\frac{\partial \mathcal{L}}{\partial \dot{q}_i} = \frac{d}{dt}p_i = 0 \implies p_i=\text{const.}
\end{equation*}
\begin{equation*}
    \mathcal{L}_{\text{effective}}=\mathcal{L}-p_i \dot{q}_i \quad \texttt{(for cyclic coordinates)}
\end{equation*}
\begin{equation*}
    p_i = \frac{\partial \mathcal{L}}{\partial \dot{q}_i \quad }
\end{equation*}
\begin{equation*}
    \frac{\partial\mathcal{L}}{\partial t}=0 \Rightarrow E=\text{const.}
\end{equation*}

\textbf{Hamiltonian Mechanics}
\begin{equation*}
    \mathcal{H} = \sum_i \dot{q}_i p_i - \mathcal{L}
\end{equation*}
\begin{equation*}
    \dot{q}_i = \frac{\partial \mathcal{H}}{\partial p_i}=\{\mathcal{H},q_i\}, \quad \dot{p}_i = -\frac{\partial \mathcal{H}}{\partial q_i}=-\{\mathcal{H},p_i\}
\end{equation*}
\[
\frac{d}{dt} f \left( p_i, q_i, t \right) = \{ \mathcal{H}, f \} + \frac{\partial f}{\partial t}
\]
\[
\frac{d}{dt} \{ f, g \} = \frac{\partial \{ f, g \}}{\partial t} + \{ \{ f, g \}, \mathcal{H} \}
\]

\textbf{Hamiltonian for Common Systems}\\
\( \mathcal{H} = \frac{p^2}{2m} + \frac{1}{2} k x^2 \quad \text{(Harmonic Oscillator)} \)\\
\( \mathcal{H} = \frac{p^2}{2m} - \frac{GMm}{r} \quad \text{(Kepler Problem)} \)

\textbf{Poisson Brackets}
\begin{equation*}
    \{ f, g \} = \sum_i \left( \frac{\partial f}{\partial q_i} \frac{\partial g}{\partial p_i} - \frac{\partial f}{\partial p_i} \frac{\partial g}{\partial q_i} \right)
\end{equation*}
\[
\{f, g\} =
\begin{vmatrix}
\frac{\partial f}{\partial q_i} & \frac{\partial f}{\partial p_i} \\
\frac{\partial g}{\partial q_i} & \frac{\partial g}{\partial p_i}
\end{vmatrix}
= \sum_i \left( \frac{\partial f}{\partial q_i} \frac{\partial g}{\partial p_i} - \frac{\partial f}{\partial p_i} \frac{\partial g}{\partial q_i} \right).
\]
\begin{align*}
 \{   q_i, q_j = \{p_i, p_j\} = 0, \quad \{p_i, q_j\} = \delta_{ij}
\end{align*}
\begin{equation*}
    \{A, B\} = -\{B, A\}
\end{equation*}
\begin{equation*}
    \{A, BC\} = B\{A, C\} + \{A, B\}C
\end{equation*}
\begin{align*}
   \{ f, \{g,h\}\} + \{g, \{h,f\}\} + \{h, \{f,g\}\} = 0
\end{align*}
\begin{align*}
  \{  f(u), g(v)\} = \{u,v\} \frac{\partial f}{\partial u} \frac{\partial g}{\partial v}
\end{align*}
\begin{align*}
    f(\vec{p}), g(\vec{q})\} = \sum \frac{\partial f}{\partial p_i} \frac{\partial g}{\partial q_i}
\end{align*}
\begin{align*}
    f(\vec{q}, \vec{p}), q_i\} = \frac{\partial f}{\partial p_i}, \quad \{ f(\vec{q}, \vec{p}), p_i\} = -\frac{\partial f}{\partial q_i}
\end{align*}
\begin{align*}
    \{L_i, L_j\} = -\sum_k \epsilon_{ijk} L_k
\end{align*}

\textbf{Constant of Motion}
\begin{equation*}
    \frac{dA}{dt} = \{ A, H \} + \frac{\partial A}{\partial t}
\end{equation*}

\textbf{Canonical Transformations Condition}

A transformation \( (q_i, p_i) \to (Q_i, P_i) \) is \textbf{canonical} if it preserves the Poisson brackets:

\[
\{Q_i, Q_j\} = 0, \quad \{P_i, P_j\} = 0, \quad \{Q_i, P_j\} = \delta_{ij}.
\]

Alternatively, it satisfies the condition:

\[
\begin{bmatrix}
\frac{\partial Q}{\partial q} & \frac{\partial Q}{\partial p} \\
\frac{\partial P}{\partial q} & \frac{\partial P}{\partial p}
\end{bmatrix}=\pm1
\]

Generating function \( F(q, Q, t) \) satisfies:

\[
p_i = \frac{\partial F}{\partial q_i}, \quad P_i = \frac{\partial F}{\partial Q_i}.
\]

\begin{equation*}
    \sum_i (p_i dq_i - P_i dQ_i) = dF
\end{equation*}
\[p\dot{q}-\mathcal{H}=P\dot{Q}-\mathcal{H}'+\frac{dF}{dt}\]
\textbf{Generating Functions}
\begin{center}
\begin{tabular}{|c|l|}
\hline
\(F_1(q,Q)\) & \( p=\dfrac{\partial F_1}{\partial q},\quad P=-\dfrac{\partial F_1}{\partial Q}\) \\[4pt] \hline
\(F_2(q,P)\) & \( p=\dfrac{\partial F_2}{\partial q},\quad Q=\dfrac{\partial F_2}{\partial P}\) \\[4pt] \hline
\(F_3(p,Q)\) & \( q=-\dfrac{\partial F_3}{\partial p},\quad P=-\dfrac{\partial F_3}{\partial Q}\) \\[4pt] \hline
\(F_4(p,P)\) & \( q=-\dfrac{\partial F_4}{\partial p},\quad Q=\dfrac{\partial F_4}{\partial P}\) \\[4pt] \hline
\end{tabular}
\end{center}

%====================================================================
% III. Oscillatory Motion and Small Oscillations
%====================================================================

\textbf{Damped Harmonic Oscillator}
\begin{equation*}
    \ddot{x} + 2\gamma\dot{x} + \omega_0^2 x = f(x)
\end{equation*}

\textbf{Fourier Transform Solution}
\begin{equation*}
    x_p = \frac{1}{\sqrt{2\pi}} \int_{-\infty}^{\infty} \frac{\hat{f}(\omega) e^{i\omega t}\, d\omega}{\omega_0^2 - \omega^2 + 2i\gamma\omega}
\end{equation*}

\textbf{Forced Oscillation}
\begin{align*}
    x_p &= \frac{f_0 \cos(\omega_d t + \varphi)}{\sqrt{(\omega_0^2 - \omega_d^2)^2 + (2\gamma\omega_d)^2}} \\
    \varphi &= \tan^{-1} \left( \frac{2\gamma\omega_d}{\omega_0^2 - \omega_d^2} \right)
\end{align*}

\textbf{Matrix Inverse}
\begin{equation*}
    A = \begin{pmatrix} a & b \\ c & d \end{pmatrix} \Rightarrow A^{-1} = \frac{1}{\det A} \begin{pmatrix} d & -b \\ -c & a \end{pmatrix}
\end{equation*}

\textbf{Small Oscillations Solution Scheme}
We assume the system has an equilibrium at $\dot{\mathbf{q}}_0$. Defining displacements \(q_i\) from equilibrium and expanding the Lagrangian to second order (first-order terms vanish), we write:
\[
\mathcal{L}_{\text{eff}} \approx \frac{1}{2} \dot{\mathbf{q}}^T M \dot{\mathbf{q}} - \frac{1}{2} \mathbf{q}^T K \mathbf{q}
\]
Now, solve the eigenvalue problem:
\[
\det \left( K - \omega^2 M \right) = 0
\]
with the solution given by
\[
\vec{q}(t)\approx\sum_i \vec{u_i}\cos(\omega_i t)
\]

\textbf{Oscillatory Motion Period}
\begin{align*}
        t &= \int_{x_0}^{x} \frac{dx}{\sqrt{\frac{2}{m} (E - V(x))}} \\
        \tau &= 2 \int_{x_1(E)}^{x_2(E)} \frac{dx}{\sqrt{\frac{2}{m} (E - V(x))}}
\end{align*}
\( x_1(E) \) and \( x_2(E) \) are the turning points defined by \( E = V(x) \).

%====================================================================
% IV. Orbital Motion and Central Force Problems
%====================================================================

\textbf{Orbital Motion Equations}

\begin{align*}
    &\textbf{Differential Equation:} \quad \frac{d^2u}{d\theta^2} + u = -\frac{m}{L^2} \frac{dV}{du} \\
    &\textbf{Orbit Equation:} \quad r(\theta) = \frac{a(1 - e^2)}{1 + e\cos(\theta - \theta_0)} \\
    &\textbf{Constants:} \quad r_0 = \frac{p_\theta^2}{\mu G m_1 m_2} = a(1 - e^2)
\end{align*}

\begin{align*}
    &\textbf{Periapsis (} r_{\min} \textbf{):} \quad r_{\min} = \frac{r_0}{1 + e} = a(1 - e) \\
    &\textbf{Apoapsis (} r_{\max} \textbf{):} \quad r_{\max} = \frac{r_0}{1 - e} = a(1 + e) \\
    &\textbf{Semi-major Axis:} \quad a = \frac{r_{\max} + r_{\min}}{2} \\
    &\textbf{Semi-minor Axis:} \quad b = a\sqrt{1 - e^2}
\end{align*}

\begin{align*}
    &\textbf{Eccentricity:} \quad e = \sqrt{1 + \frac{2E p_\theta^2}{G^2 \mu m_1^2 m_2^2}} = \sqrt{1 - \left(\frac{b}{a}\right)^2}
\end{align*}

\begin{align*}
    &\textbf{Total Energy:} \quad E = \frac{1}{2} m v^2 - \frac{G m_1 m_2}{r}
\end{align*}

\begin{align*}
    &\textbf{Radial Velocity:} \quad \dot{r} = \sqrt{\frac{2}{\mu} \left( E - V(r) - \frac{\ell^2}{2\mu r^2} \right)}
\end{align*}

\begin{align*}
    &\textbf{Angular Momentum:} \quad L = m r^2 \dot{\theta} = \text{constant}
\end{align*}

\begin{align*}
    &\textbf{Orbit Types:} \\
    &\text{Elliptical (} e < 1\text{):} \quad r(\theta) = \frac{r_0}{1 + e \cos(\theta - \theta_0)} \\
    &\text{Parabolic (} e = 1\text{):} \quad r(\theta) = \frac{p_\theta^2}{2\mu G m_1 m_2} \\
    &\text{Hyperbolic (} e > 1\text{):} \quad r(\theta) = \frac{r_0}{1 + e \cos(\theta - \theta_0)}
\end{align*}

\begin{center}
\tiny
\begin{tabular}{|>{\raggedright\arraybackslash}p{0.05\linewidth}|>{\raggedright\arraybackslash}p{0.22\linewidth}|c|} \hline 
Orbit & Condition & Key Eq. \\ \hline 

Circ & \(E=E_{\min},\,\epsilon=0\) & \(r_0=\frac{L^2}{\mu G m_1m_2}\) \\ \hline  

Ellip & \(E_{\min}<E<0,\,0<\epsilon<1\) & \(r=\frac{Gm_1m_2}{2E}\Bigl(1\pm\sqrt{1+\frac{2EL^2}{\mu (Gm_1m_2)^2}}\Bigr)\) \\ \hline  

Parab & \(E=0,\,\epsilon=1\) & \(r_{\text{par}}=\frac{L^2}{2\mu Gm_1m_2}\) \\ \hline 

Hyp & \(E>0,\,\epsilon>1\) & \(r_{\text{hyp}}=\frac{Gm_1m_2}{2E}(1+\epsilon)\) \\ \hline

\end{tabular}
\end{center}

\textbf{Kepler Two-Body Problem}
\begin{align*}
    L &= \frac{\mu}{2} \left( \dot{r}^2 + r^2 \dot{\phi}^2 \right) - V(r) \\
    \mu \ddot{r} &= -\frac{dV}{dr} + \frac{p_\theta^2}{\mu r^3} \\
    \mu &= \frac{m_1m_2}{m_1+m_2}
\end{align*}
\textit{in the COM frame:}
\begin{align*}
    L &= \frac{\mu \dot{r}^2}{2} - V(r) \\
    L &= \mathbf{r} \times \mu \dot{\mathbf{r}} = \text{const} \\
    P_{\phi} &= \frac{\partial L}{\partial \dot{\phi}} = \mu r^2 \dot{\phi} = \text{const} \\
    \mu \ddot{\mathbf{r}} &= -\frac{GM\mu}{r^2} \hat{\mathbf{r}} = -\frac{\partial V_{\text{eff}}}{\partial r}
\end{align*}
\begin{align*}
    V_{\text{eff}}(r) &= V(r) + \frac{p_\theta^2}{2\mu r^2} = - \frac{G M_1 M_2}{r} + \frac{l^2}{2\mu r^2} \\
    E &= \frac{1}{2} \mu \dot{r}^2 + V_{\text{eff}}(r)
\end{align*}
\begin{align*}
    \Delta\phi &\equiv 2\int_{r_{1}}^{r_{2}}\frac{L\, dr}{r^{2}\sqrt{2\mu(E-V_{\text{eff}}(r))}}
\end{align*}
\begin{align*}
    \Delta \theta_0 = 2 \left( \theta (r_{\max}) - \theta (r_{\min}) \right) - 2\pi
\end{align*}
 


\textbf{Binet’s Equation}
\begin{equation*}
    u''+u=-\frac{\mu }{p_\theta^{2}}\frac{1}{u^{2}}F\Bigl(\frac{1}{u}\Bigr)
\end{equation*}
For a force \( F(r) = -\frac{k}{r^2} \):
\begin{equation*}
    \frac{d^2 u}{d\theta^2} + u = \frac{k}{m h^2}
\end{equation*}
and its general solution is
\begin{equation*}
    u(\theta) = \frac{1}{r} = \frac{k}{m h^2} + C \cos(\theta - \theta_0)
\end{equation*}
For gravitation,
\begin{align*}
    u'' + u = \frac{\mu G m_1 m_2}{p_\theta^2} \quad \Rightarrow \quad r(\theta) = \frac{r_0}{1 + e \cos(\theta - \theta_0)}
\end{align*}
with
\begin{align*}
    r_0 = \frac{p_\theta^2}{\mu G m_1 m_2}
\end{align*}

%====================================================================
% V. Rigid Body Dynamics
%====================================================================

\textbf{Rigid Body Dynamics}
\begin{equation*}
    \mathbf{R}_{CM} = \frac{1}{M} \sum_i m_i \mathbf{r}_i
\end{equation*}
\begin{equation*}
    I_{ij} = \sum_k m_k \left( r_k^2 \delta_{ij} - x_{k_i} x_{k_j} \right), \int\left(r^{2}\delta_{ij}-r_{i}r_{j}\right)dm
\end{equation*}
\begin{equation*}
    I = I_C + Md^2,\quad I_{ij}^{\prime}=I_{ij}+M\left[a^{2}\delta_{ij}-a_{i}a_{j}\right]
\end{equation*}
\begin{equation*}
\begin{cases}
    I_1 \dot{\omega}_1 + (I_3 - I_2) \omega_2 \omega_3 = \tau_1 \\
    I_2 \dot{\omega}_2 + (I_1 - I_3) \omega_3 \omega_1 = \tau_2 \\
    I_3 \dot{\omega}_3 + (I_2 - I_1) \omega_1 \omega_2 = \tau_3
\end{cases}
\end{equation*}

\textbf{Moment of Inertia}
\begin{equation}
    I = \sum m_{\alpha} r_{\alpha}^2
\end{equation}
\[
I = \int_V \rho(\mathbf{r})\, r^2\, dV
\]
\begin{equation*}
    I_{xx} + I_{yy} = I_{zz} \quad \texttt{x and y in the plane of the object}
\end{equation*}
\begin{align*}
    I_\hat{n}=\hat{n}^TI\hat{n},\quad \hat{n}=(0,\sin\alpha,\cos(\alpha)
\end{align*}


\begin{center}
\renewcommand{\arraystretch}{1.1}
\setlength{\tabcolsep}{4pt}
\begin{tabular}{|l|c|}
\hline
\textbf{Object} & \(\mathbf{I}\) \\ \hline
Solid Sphere & \(\frac{2}{5}MR^2\) \\ \hline   
Hollow Sphere & \(\frac{2}{3}MR^2\) \\ \hline   
Solid Cylinder/Thin Disk & \(\frac{1}{2}MR^2\) \\ \hline   
Thin Rod (center) & \(\frac{1}{12}ML^2\) \\ \hline   
Thin Rod (end) & \(\frac{1}{3}ML^2\) \\ \hline   
Rectangular Plate (center) & \(\frac{1}{12}M(a^2+b^2)\) \\ \hline   
Rectangular Plate (edge) & \(\frac{1}{3}Ma^2\) \\ \hline   
Ring (center) & \(MR^2\) \\ \hline   
Solid Cone (sym. axis) & \(\frac{3}{10}MR^2\) \\ \hline   
Hollow Cylinder (axis) & \(M(R_1^2 + R_2^2)/2\) \\ \hline   
Solid Ellipsoid & \(\frac{1}{5}M(a^2 + b^2)\) \\ \hline   
Thin Spherical Shell & \(\frac{2}{3}MR^2\) \\ \hline  
\end{tabular}
\end{center}


\[
T=T_{C.M}+T_{Rot}=\frac{M}{2}V_{C.M}^2+\frac{I_{C.M}\omega^2}{2} \]
\[= \frac{M}{2}V_{cm}^{2}+\frac{1}{2}\vec{\omega}^{T}I\vec{\omega}
\]
\[
T_{rot}=\frac{1}{2}\sum_{i,j}\omega_{i}I_{i,j}\omega_{j}
\]
\[
\mathbf{L}_{\text{tot}} = \mathbf{L}_{\text{C.M}} + \mathbf{L}_{\text{Rot}} 
= M \mathbf{R}_{\text{C.M}} \times \mathbf{V}_{\text{C.M}} + I_{\text{C.M}} \boldsymbol{\omega}
\]

%====================================================================
% VI. Scattering
%====================================================================

\textbf{Scattering}
For scattering with impact parameter \(\rho\), scattering angle \(\theta\), asymptotic velocity \(v_\infty\), and reduced mass \(\mu\), the differential cross section is:
\begin{equation*}
    \frac{d\sigma}{d\Omega}
    \;=\;
    \frac{\rho(\theta)}{\sin \theta}
    \left|\frac{d\rho}{d\theta}\right|.
\end{equation*}
The total cross section is:
\begin{equation*}
    \sigma_{T}
    \;=\;
    2\pi \int_{0}^{\pi}
    \rho(\theta)
    \left|\frac{d\rho}{d\theta}\right|
    \,d\theta
    =2\pi\intop_{0}^{\pi}\frac{d\sigma}{d\Omega}\sin\theta\, d\theta
\end{equation*}
\[
\Delta \theta = 2 \int_{r_{\min}}^{\infty} \frac{l / r^2}{\sqrt{2\mu \left( E - V(r) - \frac{l^2}{2\mu r^2} \right)}} \, dr
\]
We identify
\[
E \;=\; \frac{1}{2}\,\mu\,v_{\infty}^{2}, \quad
\ell \;=\; \mu\,\rho\,v_{\infty},
\]
and if \(r_{\min}\) is the distance of closest approach then
\[
\rho^2=\frac{l^2}{\mu^2v_{\infty}^{2}}=\frac{l^2}{2\mu E}
\]
\[
E=V\left(r_{min}\right)+\frac{l^{2}}{2\mu r_{min}^{2}}
\]
\[
\rho^{2} = r_{\min}^{2}\Bigl(1 - \frac{V(r_{\min})}{E}\Bigr).
\]
\begin{equation*}
    \theta \approx - \frac{2b}{\mu v_\infty^2} \int_b^\infty \frac{V'(r)\, dr}{\sqrt{r^2 - b^2}}
\end{equation*}
\begin{equation*}
    b = \frac{L}{\mu v_\infty}
\end{equation*}

%====================================================================
% VII. Legendre Transform and Homogeneous Functions
%====================================================================

\textbf{Legendre Transform}
Given a function \( f(x) \), its Legendre transform is defined by:
\begin{equation*}
    \frac{df}{dx} = f'(x) = s(x)
\end{equation*}
Thus, we obtain a new function:
\begin{equation*}
    g(s) = -s\cdot x(s) + f(x(s))
\end{equation*}

\textbf{Inverse Transform}
\begin{equation*}
    V(x) = g(s(x)) + f(x)
\end{equation*}
\begin{equation*}
    -x(s) \cdot xu(x)
\end{equation*}

\textbf{Homogeneous Function of Degree \( k \)}
\( f(a x) = a^k f(x), \quad \forall a \) \quad and \quad \( E_\text{Kin}(a\vec{v})=\frac{1}{2}ma^2v^2=a^2E_\text{Kin}(\vec{v}) \)
\begin{equation*}
    V(r)=g\, r^k 
\end{equation*}

\textbf{Euler’s Theorem for Homogeneous Functions}
\begin{equation*}
    \nabla f(x) \cdot x = k f(x)
\end{equation*}
or equivalently,
\begin{equation*}
    x \cdot \frac{\partial f(x)}{\partial x} = k f(x)
\end{equation*}

\textbf{Virial Theorem}
\begin{equation*}
    2 \langle T \rangle + \langle V \rangle = 0
\end{equation*}
\begin{equation*}
    2\langle K \rangle = -\Bigl\langle \sum_j \frac{\partial \mathcal{L}}{\partial q_j} q_j \Bigr\rangle = \Bigl\langle \sum_i \vec{r}_i \frac{\partial V}{\partial \vec{r}_i} \Bigr\rangle
\end{equation*}
\[ 2\langle T\rangle = k_{\text{deg}}\langle V\rangle \]

\textbf{Scaling}
\[
x\rightarrow\alpha x,\quad t\rightarrow\beta t \Rightarrow v\rightarrow\frac{\alpha}{\beta}v
\]
\[
\tilde{\mathcal{L}}=\left(\frac{\alpha}{\beta}\right)^{2}T-\alpha^{k}V,\quad \left(\frac{\alpha}{\beta}\right)^{2}=\alpha^{k}\Rightarrow\beta=\alpha^{1-\frac{k}{2}}
\]
\begin{equation*}
    \frac{v^{\prime}}{v}=\left(\frac{l^{\prime}}{l}\right)^{\frac{1}{2} k},\quad \frac{E^{\prime}}{E}=\left(\frac{l^{\prime}}{l}\right)^k,\quad \left(\frac{l^{\prime}}{l}\right)^{1-\frac{1}{2}k}=\frac{t'}{t}
\end{equation*}

\textbf{Buckingham \(\Pi\) Theorem}
\begin{align*}
        f(x_1, \ldots, x_n) = 0 \Rightarrow f(\Pi_1, \ldots, \Pi_{n-r}) = 0
\end{align*}
\begin{equation*}
    \tau \sim d^{1 - \frac{k}{2}}, \quad v \sim \frac{d}{\tau} \sim d^{k/2}
\end{equation*}
\begin{equation*}
    E \sim m v^2 \sim d^k ,L \sim \frac{d^2}{\tau} \sim d^{1 + \frac{k}{2}}
\end{equation*}

%====================================================================
% VIII. Coordinates, Transformations and Vector Calculus
%====================================================================
\textbf{Vector and Tensor Calculus}
\begin{equation*}
    \mathbf{a} \cdot \mathbf{b} = a_i b_i, \quad    ( \mathbf{a} \times \mathbf{b} )_i = \epsilon_{ijk} a_j b_k
\end{equation*}
\begin{equation*}
    (\nabla \times \mathbf{b})_i = \epsilon_{ijk} \partial_j b_k ,\quad  (AB)_{ij} = A_{ik} B_{kj}
\end{equation*}


\textbf{Index Notation and Levi-Civita}
\begin{equation*}
    \epsilon_{ijk} \epsilon_{ilm} = \delta_{jl} \delta_{km} - \delta_{jm} \delta_{kl}
\end{equation*}
\begin{align*}
    \epsilon_{123}=1,\quad \epsilon_{213}=-1,\quad \epsilon_{221}=0
\end{align*}
\begin{equation*}
    \epsilon_{ijk}= \epsilon_{jki} = -\epsilon_{jik}
\end{equation*}
\begin{equation*}
    \hat{e}_i \times \hat{e}_j=\epsilon_{ijk} \hat{e}_k
\end{equation*}

\textbf{Complete Derivative}
\begin{equation*}
    \frac{d X}{d t}=\frac{\partial X}{\partial t}+\frac{\partial X}{\partial q} \frac{d q}{d t} 
\end{equation*}
\begin{align*}
    \mathbf{r}  &= r(\cos\varphi,\ \sin\varphi) = r \hat{\mathbf{r}}, \\[1ex]
   \dot{\mathbf{r}} &= \dot{r}\hat{r}+r\dot{\theta}\hat{\theta}, \\[1ex]
  \ddot{\mathbf{r}} &= (\ddot{r}-r\dot{\theta}^2)\hat{r}+(2\dot{r}\dot{\theta}+r\ddot{\theta})\hat{\theta}
\end{align*}

\textbf{Rotation Matrices:}
\[
\tiny R_x(\gamma) =
\begin{bmatrix}
1 & 0 & 0 \\
0 & \cos\gamma & -\sin\gamma \\
0 & \sin\gamma & \cos\gamma
\end{bmatrix},\tiny
R_y(\beta) =
\begin{bmatrix}
\cos\beta & 0 & \sin\beta \\
0 & 1 & 0 \\
-\sin\beta & 0 & \cos\beta
\end{bmatrix}
\]
\tiny
\[
R_z(\alpha) =
\begin{bmatrix}
\cos\alpha & -\sin\alpha & 0 \\
\sin\alpha & \cos\alpha & 0 \\
0 & 0 & 1
\end{bmatrix}
\]
\textbf{Combined Rotation:}  
\[
R(\alpha, \beta, \gamma) = R_z(\alpha) R_y(\beta) R_x(\gamma)
\]


\textbf{Angular Velocity Matrix}
\begin{align*}
\Omega=
    \begin{pmatrix} \omega_1 \\ \omega_2 \\ \omega_3 \end{pmatrix}
=
\begin{pmatrix} 
\dot{\phi} \sin\theta \sin\psi + \dot{\theta} \cos\psi \\[2mm]
\dot{\phi} \sin\theta \cos\psi + \dot{\theta} \sin\psi \\[2mm]
\dot{\psi} + \dot{\phi} \cos\theta 
\end{pmatrix}
\end{align*}

%====================================================================
% IX. Parametric Resonance
%====================================================================

\textbf{Parametric Resonance}
\[
\ddot{x} + \omega (t) x = 0,\quad \text{with } \omega (t)=\omega (t+T)\implies x^*(t)=x(t+T).
\]
Substituting the periodic solution space, we obtain:
\[
x_i (t + nT) = \mu_i^n x_i = \sum_j A_{ij}\tilde{x}(t)
\]
\[
x_i(t)=\mu_i^{t/T}\Pi_i(t)
\]
\[
\left|\mu_{1}\right|\neq\left|\mu_{2}\right|\implies\textbf{Parametric Resonance}
\]
with the condition:
\[
\mu_1 \mu_2 = 1 \quad \Rightarrow \quad x_2 x_1 - x_1 x_2 = \text{const}
\]
Thus, parametric resonance occurs if \(\mu_1 \neq 1\).

\textbf{Hamilton-Jacobi Equation (Single Degree of Freedom):}
\[
\frac{\partial S}{\partial t} + \mathcal{H} \left( q, \frac{\partial S}{\partial q}, t \right) = 0.
\]

\textbf{Generalization for Multiple Degrees of Freedom:}
\[
S = S \left(\{q_i\}_{i=1}^{n}, t; \{\alpha_i\}_{i=1}^{n} \right),
\quad p_i = \frac{\partial S}{\partial q_i}, \quad
\frac{\partial S}{\partial \alpha_i} = \beta_i.
\]

\textbf{Solution Procedure:}
1. Solve the Hamilton-Jacobi equation:
   \[
   \frac{\partial S}{\partial t} + \mathcal{H} \left( q, \left\{ \frac{\partial S}{\partial q_i} \right\}_{i=1}^{n}, t \right) = 0.
   \]
2. Separate variables and find \( S(q,t;\alpha) \).
3. Compute:
   \[
   \frac{\partial S}{\partial \alpha} = \beta, \quad q_i = q_i(\alpha, \beta, t).
   \]
4. Determine:
   \[
   p = \frac{\partial S}{\partial q}, \quad P = P(\alpha, \beta, t).
   \]
5. Find \( \alpha, \beta \) using initial conditions \( q(0), p(0) \)
%====================================================================
% X. Taylor Series
%====================================================================

\textbf{Taylor Expansions}
\[
(1 + x)^\alpha = 1 + \alpha x + \frac{\alpha(\alpha-1)}{2!}\,x^2 + \cdots
\]
\[
\sqrt{1+f^2} = 1 + \frac{f^2}{2} - \frac{f^4}{8} 
\]
\[
f(x) = f(a) + f'(a)\,(x - a) + \frac{f''(a)}{2!}\,(x - a)^2 + \cdots
\]
\[
\sin x = x - \frac{x^3}{3!}+ \quad \cos x = 1 - \frac{x^2}{2!}
\]

\textbf{Trigonometric Identities}
\textbf{Fundamental:}  
$1 + \tan^2\theta = \sec^2\theta$, \quad  
$1 + \cot^2\theta = \csc^2\theta$.  

\textbf{Angle Sum/Difference:}  
$\sin(A\pm B) = \sin A\cos B \pm \cos A\sin B$  
$\cos(A\pm B) = \cos A\cos B \mp \sin A\sin B$  
$\tan(A\pm B) = \frac{\tan A \pm \tan B}{1 \mp \tan A\tan B}$  

\textbf{Double Angle:}  
$\sin 2\theta = 2\sin\theta\cos\theta$, \quad  
$\cos 2\theta = \cos^2\theta - \sin^2\theta$, \quad  
$\tan 2\theta = \frac{2\tan\theta}{1 - \tan^2\theta}$.  

\textbf{Half-Angle:}  
$\sin^2\frac{\theta}{2} = \frac{1 - \cos\theta}{2}$, \quad  
$\cos^2\frac{\theta}{2} = \frac{1 + \cos\theta}{2}$, \quad  
$\tan\frac{\theta}{2} = \frac{\sin\theta}{1+\cos\theta}$.  

\textbf{Product-to-Sum:}  
$\sin A\sin B = \frac{1}{2}[\cos(A-B) - \cos(A+B)]$  
$\cos A\cos B = \frac{1}{2}[\cos(A+B) + \cos(A-B)]$  
$\sin A\cos B = \frac{1}{2}[\sin(A+B) + \sin(A-B)]$  

\textbf{Sum-to-Product:}  
$\sin A + \sin B = 2\sin\frac{A+B}{2} \cos\frac{A-B}{2}$  
$\sin A - \sin B = 2\cos\frac{A+B}{2} \sin\frac{A-B}{2}$  
$\cos A + \cos B = 2\cos\frac{A+B}{2} \cos\frac{A-B}{2}$  
$\cos A - \cos B = -2\sin\frac{A+B}{2} \sin\frac{A-B}{2}$ 
$ c^{2}=a^{2}+b^{2}-2ab\cos\gamma $
\end{multicols}



\end{document}
